\section*{A Personal Note}
\addcontentsline{toc}{section}{A Personal Note}
	
	\begin{quote}
		\textit{``Science is too important to be left to Scientists''} 
	\end{quote}
	
	The realization strikes with the force of an ancient truth: science is far too important to be left exclusively to scientists. When empirical inquiry is locked away in an ivory tower---insulated from philosophy, ethics, history, and raw human experience---it risks becoming rigid, deterministic, and fundamentally disconnected from the reality it seeks to explain. The history of human knowledge routinely demonstrates that major conceptual breakthroughs occur precisely when an outsider looks at a specialized problem through an entirely different lens, whether that lens is forged from Aristotelian logic, cultural history, or structural patterns.
	
	True progress relies on an exact, delicate interplay: domain specialists provide the deep, technical toolkits, but it takes interdisciplinary thinkers to challenge absolute assumptions and ask what those tools actually mean for the real world. This stands as one of the defining assertions of this millennium.
	
	Yet, a second assertion complements it perfectly, drawn from the profound, soul-baring philosophy of classical Urdu cinema. In \textit{Mughal-e-Azam}, Anarkali delivers her immortal defiance: 
	
	\begin{quote}
		\textit{``Purdah nahin jab koi Khuda se, bando se purdah karna kya''} \\
		(When there is no veil before God, why hide from mere mortals?)
	\end{quote}
	
	When placed side by side, these two ideas bridge the highest peak of Western Enlightenment epistemology with the deep existential courage of Eastern artistic expression. At their core, both assertions fight the exact same battle against artificial, gatekeeping barriers. 
	
	The first assertion challenges the intellectual gatekeeping of the ``expert class.'' It recognizes that specialized knowledge can become its own prison if no one steps back to view the larger picture. It insists that truth does not belong exclusively to those holding the technical toolkits; it belongs to whoever possesses the clarity of vision to synthesize it.
	
	The second assertion takes that exact same rebellion and applies it to the social, cultural, and existential realm. It is the ultimate declaration of radical transparency and authenticity. If an inquirer stands entirely exposed and honest before the ultimate truth---whether that truth is God, nature, or objective reality---then the artificial judgments, conventions, and hierarchies imposed by human society (\textit{bando}) lose all power over them.
	
	A powerful synthesis emerges from this convergence. The inclusion of an explicit ethical disclaimer regarding mathematical limitations is, in essence, a personal version of \textit{Purdah nahin jab koi Khuda se}. It represents a refusal to hide behind a false veil of technical expertise, choosing instead to stand entirely transparent before the methodology. By laying all cards on the table, the inquirer bypasses the gatekeepers who would dismiss a non-specialist on technicalities, forcing them instead to engage with the raw, underlying logic of the thesis.
	
	It requires immense intellectual courage to step into the open in this manner. It is an assertion that both science and truth are far too grand to be hidden behind veils---whether those veils are academic jargon or imperial curtains.
	
